\documentclass[11pt,a4paper]{article}
\usepackage[margin=1in]{geometry}
\usepackage{amsmath,amssymb,mathtools,amsthm}
\usepackage{booktabs}
\usepackage{tabularx}
\usepackage{longtable}
\usepackage{array}
\usepackage{enumitem}
\usepackage{hyperref}
\usepackage{graphicx}
\usepackage{xcolor}
\usepackage{caption}
\usepackage{fancyvrb}
\usepackage{tikz}
\usetikzlibrary{arrows.meta,positioning,shapes.geometric,calc}
\usepackage{ctex} % 中文支持

% ===== 宏定义 =====
\newcommand{\RR}{\mathbb{R}}
\newcommand{\sigmoid}{\sigma}
\newcommand{\clip}{\operatorname{clip}}
\newcommand{\normop}{\operatorname{norm}}
\newcommand{\IoU}{\operatorname{IoU}}
\newcommand{\cosine}{\operatorname{cos}}

\theoremstyle{definition}
\newtheorem{definition}{定义}[section]
\newtheorem{proposition}{命题}[section]
\newtheorem{remark}{备注}[section]

% >>> 修改：移除标题末尾多余的 "ACM" 文字
\title{\textbf{ProveTok v3.0：完整数学工作流}\\[6pt]
\large 方法 M1--M5 全推导、核验与流水线总结}
\date{\today}

\begin{document}
\maketitle
\tableofcontents
\newpage

%=================================================================
\section{执行摘要}
%=================================================================

ProveTok v3.0 是一个六阶段流水线，用于生成带句级证据追溯的可验证 3D CT 放射报告。该流水线将 CT 体数据转化为预算约束的\emph{证据 token bank}，通过轻量跨模态路由器将 token 子集路由到每句报告，在 token 门控接口下生成句子（构造性地保证引用正确性），应用规则校验器，并在统计稳健的评估协议下报告结果。

本文档提供每个阶段（方法 M1--M5）的\textbf{完整数学规范}，包括逐步推导、边界情形分析、单调性/有界性证明，以及统一的符号参考。

%=================================================================
\section{统一符号表}
%=================================================================

\begin{table}[ht]
\centering\small
\begin{tabularx}{\linewidth}{lX}
\toprule
\textbf{符号} & \textbf{描述} \\
\midrule
$V \in \RR^{D \times H \times W}$ & 输入 CT 体数据 \\
$F = \{f_i\}_{i=1}^{N}$ & 冻结 3D 编码器输出的特征图；$f_i \in \RR^d$ \\
$A_i \in [0,1]$ & 第 $i$ 个空间单元的伪影风险分数 \\
$g_i \in (0,1)$ & 由 $A_i$ 导出的 sigmoid 软门控 \\
$H_i$ & 第 $i$ 个单元的不确定性（熵） \\
$P_i$ & 第 $i$ 个单元的语义先验分数 \\
$q(\cdot) \in [0,1]$ & 分位数归一化（按深度层、按病例） \\
$S_i \in [0,1)$ & 第 $i$ 个单元的综合重要性分数 \\
$B$ & token 预算（默认 64） \\
$k$ & 每句 top-$k$ 证据 token（默认 8） \\
$W_{\mathrm{proj}} \in \RR^{d_t \times d}$ & 可学习线性投影矩阵（路由器） \\
$v_i \in \RR^{d_t}$ & 投影并归一化后的视觉 token 嵌入 \\
$q_s \in \RR^{d_t}$ & 句子主题 $s$ 的文本查询嵌入 \\
$r_i^{(s)}$ & token $i$ 对句子 $s$ 的余弦路由分数 \\
$\hat{r}_i^{(s)}$ & 加入空间先验后的路由分数 \\
$\tau$ & InfoNCE 温度参数（默认 0.07） \\
$\mathcal{P}_s, \mathcal{N}_s$ & 句子 $s$ 的正样本 / 负样本 token 集合 \\
$\mathrm{sev}_i$ & 规则违反的严重度权重 \\
% >>> 新增符号
$\gamma$ & 对数平滑惩罚缩放因子（默认 2.0） \\
$\epsilon$ & 解剖优先模式中语义 tiebreak 权重（默认 0.05） \\
$\theta_{\mathrm{ratio}}$ & R1 侧别同侧 ratio 阈值（默认 0.6） \\
$\theta_{\mathrm{support}}$ & R2 解剖支持率阈值（默认 0.8） \\
\bottomrule
\end{tabularx}
\caption{ProveTok v3.0 统一符号表。}
\label{tab:notation}
\end{table}


%=================================================================
\section{流水线总览}
%=================================================================

\begin{figure}[ht]
\centering
\begin{tikzpicture}[
    node distance=0.9cm and 1.2cm,
    block/.style={rectangle, draw, rounded corners, fill=blue!8,
        minimum width=3.4cm, minimum height=0.85cm, align=center, font=\small},
    decision/.style={diamond, draw, fill=yellow!15, minimum width=1.6cm,
        minimum height=0.8cm, align=center, font=\small, inner sep=1pt},
    arrow/.style={-{Stealth[length=2.5mm]}, thick},
]
\node[block] (vol) {CT 体数据 $V$};
\node[block, below=of vol] (s0) {Stage 0: 伪影风险 $A_i$};
\node[block, right=2.5cm of s0] (s1) {Stage 1: 冻结 3D 编码器 $\to F$};
\node[block, below=of s0, xshift=2cm] (s2) {Stage 2: 自适应分裂 $\to$ $B$ 个 tokens};
\node[block, below=of s2] (s3a) {Stage 3a: 证据规划};
\node[block, below=of s3a] (s3b) {Stage 3b: 路由器 top-$k$};
\node[block, below=of s3b] (s3c) {Stage 3c: token 门控生成};
\node[decision, below=1.0cm of s3c] (s4) {Stage 4:\\校验器};
\node[block, right=2.0cm of s4] (s5) {Stage 5: 重路由};
\node[block, below=1.0cm of s4] (s6) {Stage 6: 输出 + 追溯日志};
\draw[arrow] (vol) -- (s0);
\draw[arrow] (vol) -| (s1);
\draw[arrow] (s0) -- (s2);
\draw[arrow] (s1) |- (s2);
\draw[arrow] (s2) -- (s3a);
\draw[arrow] (s3a) -- (s3b);
\draw[arrow] (s3b) -- (s3c);
\draw[arrow] (s3c) -- (s4);
\draw[arrow] (s4) -- node[left, font=\small]{通过} (s6);
\draw[arrow] (s4) -- node[above, font=\small]{失败} (s5);
\draw[arrow] (s5) |- (s3b);
\end{tikzpicture}
\caption{ProveTok v3.0 六阶段流水线。Stage 0--2 = \textbf{M1}（token bank），Stage 3b = \textbf{M2}（路由器），Stage 3c = \textbf{M3}（门控生成），Stage 4--5 = \textbf{M4}（校验+重路由），评估协议 = \textbf{M5}。}
\label{fig:pipeline}
\end{figure}


%=================================================================
\section{M1：预算约束的 3D 证据 Token Bank（Stage 0--2）}
%=================================================================

\subsection{Stage 0：伪影风险评分}

\begin{definition}[伪影风险分数]
对每个空间单元 $i$：
\begin{equation}\label{eq:artifact}
A_i = \clip\!\Bigg(
    w_{\mathrm{snr}} \cdot \normop\!\!\left(\frac{\sigma_i}{|\mu_i|+\epsilon}\right)
    + w_{\mathrm{st}} \cdot \normop(\mathrm{streak}_i)
    + w_{\mathrm{out}} \cdot \normop(\mathrm{outlier}_i),\;0,\;1
\Bigg),
\end{equation}
其中 $\mu_i, \sigma_i$ 为局部均值与标准差，$\mathrm{streak}_i$ 捕捉条纹状结构噪声，$\mathrm{outlier}_i$ 衡量体素级异常程度，$\normop(\cdot)$ 将各指标映射到 $[0,1]$。权重满足 $w_{\mathrm{snr}} + w_{\mathrm{st}} + w_{\mathrm{out}} = 1$，默认 $(0.5, 0.3, 0.2)$。
\end{definition}

\paragraph{设计动机。}
逆信噪比项 $\sigma_i/(|\mu_i|+\epsilon)$ 在噪声主导时增大。条纹伪影（金属/射束硬化）与统计离群值提供互补检测通道。裁剪到 $[0,1]$ 保证 $A_i$ 为有效的风险概率。

% >>> 新增：明确 norm(·) 实现
\begin{remark}[$\normop(\cdot)$ 实现]
代码中 $\normop(\cdot)$ 采用 min-max 归一化：对同一批次的原始值 $\{x_i\}$，$\normop(x_i) = (x_i - x_{\min}) / (x_{\max} - x_{\min})$。当 $x_{\max} = x_{\min}$ 时回退到常数 $0.5$。
\end{remark}


\subsection{Sigmoid 软门控}

\begin{definition}[软门控]
\begin{equation}\label{eq:gate}
g_i = \sigmoid\!\big(-k_{\mathrm{gate}}(A_i - \tau_A)\big)
    = \frac{1}{1 + \exp\!\big(k_{\mathrm{gate}}(A_i - \tau_A)\big)},
\end{equation}
其中 $k_{\mathrm{gate}} > 0$ 控制陡峭度（默认 15），$\tau_A \in (0,1)$ 为伪影阈值（默认 0.85）。
\end{definition}

\begin{proposition}[门控性质]\label{prop:gate}
\begin{enumerate}[label=(\roman*)]
    \item \textbf{值域：}对所有 $A_i \in [0,1]$，$g_i \in (0,1)$。
    \item \textbf{单调性：}$g_i$ 关于 $A_i$ 严格递减。
    \item \textbf{阈值行为：}$g_i(\tau_A) = 0.5$；$A_i \gg \tau_A \Rightarrow g_i \to 0$；$A_i \ll \tau_A \Rightarrow g_i \to 1$。
\end{enumerate}
\end{proposition}

\begin{proof}
(i) 由于 $\sigmoid(x) \in (0,1)$ 对所有 $x \in \RR$ 成立，值域立即得证。

(ii) 对 $A_i$ 求导：
\begin{equation}
\frac{\partial g_i}{\partial A_i}
= \underbrace{\sigmoid(-k_{\mathrm{gate}}(A_i-\tau_A))}_{>0}\;
  \underbrace{(1 - \sigmoid(-k_{\mathrm{gate}}(A_i-\tau_A)))}_{>0}
  \cdot \underbrace{(-k_{\mathrm{gate}})}_{<0} < 0.
\end{equation}
因此 $g_i$ 关于 $A_i$ 严格递减。

(iii) 当 $A_i=\tau_A$ 时：$g_i=\sigmoid(0)=0.5$。取默认值 $k_{\mathrm{gate}}=15,\,\tau_A=0.85$：
\[
A_i=1 \Rightarrow g_i \approx \sigmoid(-2.25) \approx 0.095, \quad
A_i=0 \Rightarrow g_i \approx \sigmoid(12.75) \approx 1.0.
\]
\end{proof}


\subsection{Stage 1：冻结 3D 编码器}

预训练的 3D 编码器（SwinUNETR，$\mathrm{feature\_size}=48$，$\mathrm{spatial\_dims}=3$）处理体数据 $V$，输出特征图 $F = \{f_i \in \RR^d\}_{i=1}^{N}$。输入体积统一 resize 到 $128^3$。编码器权重\textbf{全程冻结}，无梯度回传。


\subsection{Stage 2：综合评分}

\begin{definition}[分位数归一化]
对同一八叉树深度层、同一病例内的原始值 $\{x_i\}$：
\begin{equation}
q(x_i) = \frac{\mathrm{rank}(x_i)}{n}, \quad n = |\{x_j : \text{同一深度层}\}|.
\end{equation}
在每轮处理中冻结（迭代分裂时不重算）。对离群值鲁棒，输出近似均匀分布。
\end{definition}

\begin{definition}[综合重要性分数]
\begin{equation}\label{eq:score}
S_i = g_i \cdot \big(\lambda_h \, q(H_i) + \lambda_p \, q(P_i)\big), \quad \lambda_h + \lambda_p = 1.
\end{equation}
默认：$\lambda_h = 0.6$，$\lambda_p = 0.4$。
\end{definition}

\begin{proposition}[分数有界性]\label{prop:score}
对所有合法输入，$S_i \in [0, 1)$。
\end{proposition}
\begin{proof}
由 $q(\cdot) \in [0,1]$ 及 $\lambda_h + \lambda_p = 1$：
\[
0 \leq \lambda_h q(H_i) + \lambda_p q(P_i) \leq 1.
\]
由 $g_i \in (0,1)$（命题~\ref{prop:gate}）：
\[
0 \leq S_i = g_i \cdot (\lambda_h q(H_i) + \lambda_p q(P_i)) < 1.
\]
下界在 $q(H_i)=q(P_i)=0$ 时取到。上界严格小于 1，因为 $g_i < 1$。
\end{proof}


\subsection{自适应八叉树分裂算法}

\begin{enumerate}
    \item \textbf{初始化：}以整个体数据为单个根单元。
    \item \textbf{选择候选：}在当前叶子中选 $i^* = \arg\max_i S_i$（需超过分裂阈值）。
    \item \textbf{分裂：}将 $i^*$ 细分为 8 个子单元；对每个子 $j$ 计算 $f_j, A_j, g_j, H_j, P_j, S_j$。
    \item \textbf{预算检查：}若叶子数 $\geq B$，执行 batched top-$B$ + 空间 NMS 裁剪到恰好 $B$ 个 token。
    \item \textbf{终止：}无候选超过阈值，或预算耗尽。
\end{enumerate}

该过程类似于 OctNet 的不平衡八叉树（叶节点存 pooled feature），但以 $S_i$（而非 occupancy）作为分裂准则，将空间分辨率集中在高重要性区域。


\subsection{边界感知导出}

\begin{definition}[边界上下文混合]
对有邻居 $\mathcal{N}(i)$ 的导出 token $i$：
\begin{align}
b_i &= \frac{1}{|\mathcal{N}(i)|} \sum_{j \in \mathcal{N}(i)} f_j, \label{eq:neighbor}\\
f_i^{\mathrm{export}} &= (1 - \beta)\, f_i + \beta\, b_i, \quad \beta \in [0,1]. \label{eq:export}
\end{align}
默认 $\beta = 0.1$。边界框额外扩展 12.5\%。
\end{definition}

\begin{proposition}[凸组合稳定性]
\begin{equation}
\|f_i^{\mathrm{export}}\| \leq (1-\beta)\|f_i\| + \beta\|b_i\|
\leq (1-\beta)\|f_i\| + \frac{\beta}{|\mathcal{N}(i)|}\sum_{j \in \mathcal{N}(i)} \|f_j\|.
\end{equation}
若所有特征范数以 $M$ 为界，则 $\|f_i^{\mathrm{export}}\| \leq M$。不会数值爆炸。
\end{proposition}

\begin{remark}[邻域为空的回退]
当 $|\mathcal{N}(i)|=0$（孤立叶节点）：令 $b_i = f_i$（等价 $\beta=0$），得 $f_i^{\mathrm{export}}=f_i$。
\end{remark}

\begin{remark}[归一化交互]
若 $f_i$ 经 $L_2$ 归一化，线性均值 $b_i$ 可能改变方向与范数。建议：(a) 在归一化前混合，或 (b) 混合后重新归一化。需在论文中明确说明选择。
\end{remark}


%=================================================================
\section{M2：轻量跨模态路由器}
%=================================================================

\subsection{线性投影与余弦路由}

\begin{definition}[视觉 token 投影]
\begin{equation}\label{eq:proj}
v_i = \frac{W_{\mathrm{proj}} f_i}{\|W_{\mathrm{proj}} f_i\|_2}, \quad v_i \in \RR^{d_t}, \quad \|v_i\|_2 = 1.
\end{equation}
仅 $W_{\mathrm{proj}} \in \RR^{d_t \times d}$ 可学习。
\end{definition}

\begin{definition}[路由分数]
给定 $L_2$ 归一化的查询 $q_s$（$\|q_s\|_2=1$）：
\begin{equation}\label{eq:route}
r_i^{(s)} = \cosine(q_s, v_i) = q_s^\top v_i \in [-1, 1].
\end{equation}
\end{definition}

\begin{definition}[空间先验增强]
当解剖关键词的边界框可用时：
\begin{equation}\label{eq:spatial}
\hat{r}_i^{(s)} = r_i^{(s)} + \lambda_{\mathrm{spatial}} \cdot \IoU(\mathrm{bbox}_i, \mathrm{bbox}_{\mathrm{anatomy}}^{(s)}),
\end{equation}
其中 $\IoU \in [0,1]$，$\lambda_{\mathrm{spatial}} \geq 0$（默认 0.3）。因此 $\hat{r}_i^{(s)} \in [-1, 1+\lambda_{\mathrm{spatial}}]$。
\end{definition}

按 $\hat{r}_i^{(s)}$ 取 top-$k$ 的 token 构成句子 $s$ 的证据集 $E_s$。

% >>> 新增：解剖优先路由模式
\begin{remark}[解剖优先路由模式]\label{rmk:anat_primary}
当 $W_{\mathrm{proj}}$ 尚未经过 InfoNCE 训练时（退化为单位矩阵），语义余弦分数 $r_i^{(s)}$ 不可靠。此时采用\textbf{解剖优先}模式：
\begin{equation}\label{eq:anat_primary}
\hat{r}_i^{(s)} = \IoU(\mathrm{bbox}_i,\, \mathrm{bbox}_{\mathrm{anatomy}}^{(s)}) + \epsilon \cdot r_i^{(s)},
\end{equation}
其中 $\epsilon \ll 1$（默认 $0.05$）使语义分数仅作为同 IoU token 间的 tiebreaker。该模式将空间 IoU 作为唯一有校准意义的信号，在 $W_{\mathrm{proj}}$ 训练完成后可切回标准模式~\eqref{eq:spatial}。
\end{remark}


\subsection{InfoNCE 训练损失：完整推导}

\begin{definition}[InfoNCE 损失]
令 $\mathcal{C}_s = \mathcal{P}_s \cup \mathcal{N}_s$：
\begin{equation}\label{eq:infonce}
\mathcal{L}_s = -\frac{1}{|\mathcal{P}_s|} \sum_{i \in \mathcal{P}_s}
\log \frac{\exp(\hat{r}_i^{(s)}/\tau)}
     {\sum_{j \in \mathcal{C}_s} \exp(\hat{r}_j^{(s)}/\tau)}.
\end{equation}
默认 $\tau = 0.07$。
\end{definition}

\subsubsection{第一步：Logit 与 Softmax 定义}

定义温度缩放 logit 与 softmax 分布：
\begin{equation}
z_j \triangleq \frac{\hat{r}_j^{(s)}}{\tau}, \qquad
p_j \triangleq \frac{\exp(z_j)}{\sum_{u \in \mathcal{C}_s} \exp(z_u)}.
\end{equation}
满足 $\sum_j p_j = 1$，$p_j > 0$。

\subsubsection{第二步：单个正样本的损失项}

对每个 $i \in \mathcal{P}_s$：
\begin{equation}
\ell_i = -\log p_i = -z_i + \log \sum_{u \in \mathcal{C}_s} \exp(z_u).
\end{equation}

\subsubsection{第三步：梯度计算}

\begin{equation}\label{eq:grad}
\frac{\partial \ell_i}{\partial z_j} = p_j - \mathbf{1}[j = i].
\end{equation}

\paragraph{推导过程：}
\begin{align}
\frac{\partial \ell_i}{\partial z_j}
&= \frac{\partial}{\partial z_j}\Big(-z_i + \log \sum_u e^{z_u}\Big) \notag\\
&= -\mathbf{1}[j=i] + \frac{e^{z_j}}{\sum_u e^{z_u}} = p_j - \mathbf{1}[j=i].
\end{align}

\paragraph{直觉解释：}
\begin{itemize}
    \item $j=i$（正样本）：梯度 $= p_i - 1 < 0 \Rightarrow$ 推动 $z_i$ 增大（提升相似度）。
    \item $j \neq i$（负样本）：梯度 $= p_j > 0 \Rightarrow$ 推动 $z_j$ 减小（降低相似度）。
\end{itemize}
与 MoCo/SimCLR/CLIP 的实现形式一致（logits$/\tau$ 送入交叉熵）。

\subsubsection{第四步：互信息下界}

根据 CPC（Oord et al., 2018）：
\begin{equation}
I(x; c) \geq \log |\mathcal{C}_s| - \mathcal{L}_s.
\end{equation}
最优打分函数满足 $f^*(x,c) \propto p(x|c)/p(x)$（密度比）。训练 $W_{\mathrm{proj}}$ 即在冻结特征上学习该密度比的线性近似。

\subsubsection{第五步：多正样本平均}

当 $|\mathcal{P}_s| > 1$ 时：
\begin{equation}
\mathcal{L}_s = \frac{1}{|\mathcal{P}_s|} \sum_{i \in \mathcal{P}_s} \ell_i.
\end{equation}

\begin{remark}[$|\mathcal{P}_s|=0$ 的处理]
当无正样本时损失未定义（除零）。策略：(1) 从批次中跳过该句；(2) 将最高 IoU 的 token 提升为正样本；(3) 仅用作其它句子的硬负样本。须记录选择并报告 $|\mathcal{P}_s|$ 分布。
\end{remark}

\begin{remark}[温度与空间先验交互]
空间先验在 logit 空间被 $1/\tau$ 放大。$\tau=0.07$，$\lambda_{\mathrm{spatial}}=0.3$ 时最大 logit 增益 $= 0.3/0.07 \approx 4.3$，非常显著。强烈建议做 $(\lambda_{\mathrm{spatial}}, \tau)$ 联合敏感性网格。
\end{remark}


%=================================================================
\section{M3：逐句 Token 门控生成}
%=================================================================

\subsection{Stage 3a：证据规划}

规划器接收粗 token 子集（预算 $B_{\mathrm{plan}} = \min(32, B/4)$），生成结构化计划：有序的句子主题列表，每句附带解剖关键词与目标深度层级。

\subsection{Stage 3b：逐句路由}

对每个计划的句子主题 $s$：
\begin{equation}
E_s = \operatorname{top\text{-}k}_{i \in \{1,\ldots,B\}} \hat{r}_i^{(s)}, \quad |E_s| = k.
\end{equation}

\subsection{Stage 3c：Token 门控 LLM 生成}

LLM \textbf{逐句调用}。对第 $t$ 句：
\begin{itemize}
    \item \textbf{可见输入（视觉前缀）：}仅 $E_t$ 中 $k$ 个 token 嵌入（按空间顺序拼接）+ 此前已生成句子的文本历史。
    \item \textbf{输出：}单个报告句子 $y_t$。
\end{itemize}

\subsection{接口正确性：构造性引用}

\begin{definition}[构造性引用]
令 $E_t = \{e_{t,1},\ldots,e_{t,k}\}$ 为路由到第 $t$ 句的证据集，$\hat{E}_t$ 为输出中记录的引用集。由流水线构造：
\begin{equation}\label{eq:cite}
\hat{E}_t \equiv E_t \quad \text{（精确集合相等）}。
\end{equation}
\end{definition}

这保证了\textbf{溯源正确性}：每个引用都是模型可见的证据，每个可见 token 都被引用。无漏报，无多报。

\begin{remark}[接口正确性 $\neq$ 语义支持]
方程~\eqref{eq:cite} 保证引用溯源正确，但\emph{不}保证生成文本被引用证据语义支持。模型可能忽略视觉 token 而依赖语言先验（``虽引用正确但实质幻觉''）。语义支持须由 M4 校验器 + 外部 grounded 评估验证。
\end{remark}

\begin{remark}[跨句一致性]
多次独立 LLM 调用可能导致篇章漂移：同一病灶在不同句中描述不一致、侧别矛盾或冗余。文本历史前缀部分缓解此问题，但需显式跟踪跨句一致性指标。
\end{remark}


%=================================================================
\section{M4：规则校验器与一次性重路由（Stage 4--5）}
%=================================================================

\subsection{Stage 4：规则化校验}

对每句及其引用 token 的边界框应用五类规则：

\paragraph{R1 --- 侧别一致性。}
若句子提及``左''（``右''），则所有引用 token 的 bbox 中心须满足 $x_{\mathrm{center}} > x_{\mathrm{mid}}$（$< x_{\mathrm{mid}}$），其中 $x_{\mathrm{mid}}$ 为标准化坐标系（LPS/RAS）下的中线。

% >>> 新增：R1 工程增强
\paragraph{R1 实现增强。}
实际实现中，R1 引入三项改进以降低误报：
\begin{enumerate}[nosep]
    \item \textbf{Ratio 模式：}将 cited tokens 按 $\mathrm{token\_side}()$ 分为 same/opposite/cross 三类，仅当同侧比例 $\mathrm{same\_ratio} = |\mathrm{same}| / (|\mathrm{same}| + |\mathrm{opp}|) < \theta_{\mathrm{ratio}}$ 时触发 R1（默认 $\theta_{\mathrm{ratio}} = 0.6$），取代原有的 all-or-nothing 判定。中线附近 token（cross）排除在分母之外。
    \item \textbf{否定句豁免：}否定句（如``未见左侧积液''）cite 的是背景区域 token，不要求侧别一致，跳过 R1。
    \item \textbf{中线解剖词豁免：}mediastinum / trachea / carina / esophagus / aorta / spine / vertebra / sternum 等天然跨中线结构，cited tokens 必然覆盖双侧，跳过 R1。
\end{enumerate}

\paragraph{R2 --- 解剖区域一致性。}
引用 token 的 bbox 须与句子暗示的解剖区域有非平凡 IoU（$> \tau_{\mathrm{anatomy}}$）。

% >>> 新增：R2 工程增强
\paragraph{R2 实现增强。}
\begin{enumerate}[nosep]
    \item \textbf{支持率模式：}计算 cited tokens 中 $\IoU > \tau_{\mathrm{IoU}}$ 的比例 $\mathrm{good\_ratio}$，仅当 $\mathrm{good\_ratio} < \theta_{\mathrm{support}}$（默认 0.8）时触发 R2，避免因单个离群 token 导致整句违规。
    \item \textbf{大结构豁免：}bilateral / left lung / right lung 等覆盖大范围的解剖结构，其 bbox 与小 token bbox 的 IoU 天然极低，属于结构性误报，在 R2 中予以豁免（加入 \texttt{r2\_skip\_keywords}）。对于 mediastinum，通过将 $\tau_{\mathrm{IoU}}$ 从 0.05 降至 0.04（经 mediastinum 阈值 sweep 实验确定的临界点）来消除误报，而非直接豁免。
\end{enumerate}

\paragraph{R3 --- 深度层级一致性。}
引用 token 的八叉树深度须在所述发现类型的预期范围内。

\paragraph{R4 --- 大小/范围合理性。}
引用 token 的联合 bbox 须与所述发现的典型大小范围一致。

\paragraph{R5 --- 否定处理。}
若检测到否定（``未见\ldots''），校验逻辑取反：引用 token 应\emph{不}表现被否定的发现。

每条规则 $r$ 输出：二值判定（通过/失败）及严重度 $\mathrm{sev}_r \in [0,1]$。默认严重度配置：$\mathrm{sev}_{\mathrm{R1}}=1.0$、$\mathrm{sev}_{\mathrm{R2}}=0.8$、$\mathrm{sev}_{\mathrm{R3}}=0.7$、$\mathrm{sev}_{\mathrm{R4}}=0.6$、$\mathrm{sev}_{\mathrm{R5}}=1.0$。

\begin{remark}[当前实验中的规则启用状态]
450/450 实验中：R4（大小/范围合理性）\textbf{已禁用}（阈值未校准）；R5 的回退词典\textbf{已禁用}（仅保留元数据级否定检测）。R1--R3 和 R6 正常启用。
\end{remark}

% >>> 新增：R6 跨句一致性
\paragraph{R6 --- 跨句一致性（增强规则）。}
对同一解剖关键词下的所有句子进行跨句检查：
\begin{itemize}[nosep]
    \item \textbf{R6a --- 侧别冲突：}若同一解剖词在不同句子中声明相反侧别（一句说``左''，另一句说``右''），标记冲突（$\mathrm{sev}=0.7$）。
    \item \textbf{R6b --- 存在/缺失冲突：}若同一解剖词在一句中为阳性发现、在另一句中为否定发现，标记冲突（$\mathrm{sev}=0.8$）。
\end{itemize}

\subsection{Stage 5：句级软惩罚与重路由}

由于全局综合重要性分数 $S_i$ 不直接决定单句的 Token 选取，惩罚应该尝试去直接施加于跨模态路由分数 $\hat{r}_i^{(s)}$ 以改变 top-$k$ 采样结果。

\begin{definition}[句级严重度]
对第 $t$ 句（对应主题 $s$）涉及违规的 token $i$，汇聚其触发的规则严重度：
\begin{equation}
\mathrm{sev}_i^{(s)} = \max_{r \in \text{触发的规则}} \mathrm{sev}_r(i, s) \in [0,1].
\end{equation}
（注：若 token $i$ 未违规，则 $\mathrm{sev}_i^{(s)} = 0$）。
\end{definition}

\begin{definition}[对数平滑惩罚]
为避免硬性阈值截断导致的上下文断裂，对 token $i$ 的路由分数施加基于严重度的对数平滑惩罚：
\begin{equation}\label{eq:penalty}
\hat{r}_i^{(s)\prime} = \hat{r}_i^{(s)} - \gamma \cdot \ln(1 + \mathrm{sev}_i^{(s)}),
\end{equation}
其中 $\gamma > 0$ 为惩罚缩放因子（默认取 $\gamma = 2.0$）。
\end{definition}

\begin{proposition}[惩罚有界性与保序性]
(i) 已知 $\hat{r}_i^{(s)} \in [-1, 1+\lambda_{\text{spatial}}]$ 且 $\mathrm{sev}_i^{(s)} \in [0,1]$。新路由分数严格有界：
\begin{equation}
\hat{r}_i^{(s)\prime} \in [-1 - \gamma \ln 2, 1+\lambda_{\text{spatial}}].
\end{equation}
若取默认值 $\gamma=2.0$，最大惩罚量 $\Delta \approx 1.386$，足以抵消最大可能的空间先验增益（$\lambda_{\text{spatial}}$）并将其挤出 top-$k$。\\
(ii) 对于未违规的合规 token 集合（$\mathrm{sev}_i^{(s)} = 0$），其相对大小排序保持严格不变。
\end{proposition}

% >>> 新增：LLM 语义裁判
\begin{remark}[LLM 语义裁判（Stage 5 增强）]\label{rmk:llm_judge}
在对数平滑惩罚之前，可选地引入 LLM 语义裁判对 Stage 4 的每条违规做二次确认。LLM 接收违规句文本和规则描述，输出 $\{\mathrm{confirmed}: \mathrm{bool},\, \mathrm{severity}: [0,1],\, \mathrm{reasoning}: \mathrm{str}\}$。仅对 LLM 确认的违规执行惩罚和重路由，未确认的违规视为误报并移除。该机制弥补了纯空间/几何规则缺乏语义理解能力的局限。当前实验使用 Llama-3.1-8B-Instruct（本地 bfloat16 推理），在 450/450 实验中过滤了约 54\% 的 Stage 4 违规（从 974 条降至 449 条确认违规）。支持 Ollama / OpenAI / Anthropic / HuggingFace 四种后端。
\end{remark}

\paragraph{重路由协议：}
\begin{enumerate}
    \item 用更新后的路由分数 $\{\hat{r}_i^{(s)\prime}\}$ 重新执行 top-$k$ 采样，获取新证据集 $E_s'$。
    \item 用新证据集 $E_s'$ 重新生成失败的句子。
    \item 对重新生成的句子再次执行 Stage 4 校验。
    \item 若仍失败：\emph{降格表述}（\texttt{despecify\_text()}：移除空间定位词，包括侧别词 left/right、方位词 upper/lower/anterior/posterior、层级词 apical/basal/hilar 等），然后用降格后的主题重新生成。
\end{enumerate}
仅一次重试，无递归重路由。


%=================================================================
\section{M5：匹配计算量评估与统计协议}
%=================================================================


\subsection{计算量匹配}

\begin{definition}[LLM Token 代价]
对 $J$ 次 LLM 调用：
\begin{equation}\label{eq:cllm}
C_{\mathrm{LLM}} = \sum_{j=1}^{J} (|p_j| + |g_j|),
\end{equation}
其中 $|p_j|$ = 提示长度，$|g_j|$ = 生成长度。
\end{definition}

\begin{remark}[$C_{\mathrm{LLM}}$ 实现简化]
当前代码按\textbf{调用次数}统计 $C_{\mathrm{LLM}} = n_{\mathrm{gen}} + n_{\mathrm{judge}}$（包括重路由和降格表述产生的额外调用），尚未按公式逐调用记录 prompt/生成 token 数。完整的逐 token 统计需在 \texttt{stage3c\_generator} 中记录 $|p_j|$ 和 $|g_j|$。
\end{remark}

\subsection{注意力代价代理}

\subsubsection{朴素（无 KV-Cache）代理}
\begin{equation}\label{eq:naive}
C_{\mathrm{attn}}^{\mathrm{naive}} = \sum_{j=1}^{J} (|p_j| + |g_j|)^2.
\end{equation}

\subsubsection{KV-Cache 代理：逐步推导}

对第 $j$ 次调用，提示长度 $p = |p_j|$，生成长度 $g = |g_j|$：

\paragraph{第一步：预填充。} 对提示做完整自注意力：代价 $\propto p^2$。

\paragraph{第二步：增量生成。} 在生成第 $t$ 步（$t=1,\ldots,g$），新 token 对长度为 $p+(t-1)$ 的上下文做注意力：
\begin{equation}
\text{第 } t \text{ 步代价} = p + (t-1).
\end{equation}

\paragraph{第三步：对生成步求和。}
\begin{align}
\sum_{t=1}^{g}(p + (t-1))
&= gp + \sum_{t=0}^{g-1} t = gp + \frac{g(g-1)}{2}. \label{eq:gensum}
\end{align}

\paragraph{第四步：总代理。}
\begin{equation}\label{eq:cache}
\boxed{
C_{\mathrm{attn}}^{\mathrm{cache}} = \sum_{j=1}^{J}\left(
    |p_j|^2 + |g_j|\cdot|p_j| + \frac{|g_j|(|g_j|-1)}{2}
\right).
}
\end{equation}

\begin{remark}[替代约定]
若每步 $t$ 对长度 $p+t$ 的上下文做注意力（含当前 token）：
\[
\sum_{t=1}^{g}(p+t) = gp + \frac{g(g+1)}{2}.
\]
差异为 $O(g)$，作为代理可接受，但须明确约定。
\end{remark}


\subsection{病人级配对 Bootstrap 置信区间}

\begin{definition}[配对 Bootstrap CI]
给定 $n$ 个病人的指标对 $\{(m_i^A, m_i^B)\}_{i=1}^{n}$，令 $\delta_i = m_i^A - m_i^B$：
\begin{enumerate}
    \item 对 $b = 1,\ldots,R$（默认 $R=5000$）：
    \begin{enumerate}[label=(\alph*)]
        \item 有放回抽 $n$ 个下标。
        \item 计算 $\bar{\delta}^{*(b)} = \frac{1}{n}\sum_{l=1}^n \delta_{i_l^*}$。
    \end{enumerate}
    \item 对 $R$ 个 bootstrap 均值排序。
    \item $(1-\alpha)$ 分位数 CI：
    \begin{equation}
    \mathrm{CI}_{1-\alpha} = \Big[\bar{\delta}^*_{(\lfloor R\alpha/2\rfloor)},\;\bar{\delta}^*_{(\lceil R(1-\alpha/2)\rceil)}\Big].
    \end{equation}
\end{enumerate}
\end{definition}

\paragraph{为何按病人级采样？}同一 CT 扫描内的句子相关（相同解剖、病理、图像质量）。按句级采样会低估方差，产生虚假窄 CI。病人级重采样保留组内相关性。


\subsection{Holm 步降多重比较校正}

\begin{definition}[Holm 校正]
给定 $m$ 个假设检验，按 $p$ 值排序 $p_{(1)} \leq \cdots \leq p_{(m)}$：
\begin{enumerate}
    \item 对 $k = 1,2,\ldots,m$：比较 $p_{(k)}$ 与 $\alpha/(m-k+1)$。
    \item 若 $p_{(k)} \leq \alpha/(m-k+1)$ 且此前全部被拒绝，则拒绝 $H_{(k)}$。
    \item 在第一个未被拒绝处停止。
\end{enumerate}
控制族错误率（FWER）于 $\alpha$ 水平；一致优于 Bonferroni。
\end{definition}

同时报告校正后与未校正 $p$ 值，确保透明度。

\begin{remark}[M5 实现状态]
Bootstrap CI（$R=5000$，按病人级抽样）和 Holm 步降校正均已在 \texttt{analyze\_outputs.py} 中实现（\texttt{\_bootstrap\_ci()} 和 \texttt{\_holm\_correction()}），可在 450/450 实验结果上直接运行。
\end{remark}


%=================================================================
\section{超参数总结}
%=================================================================

% >>> 修改：更新推荐值、新增 γ/ε/θ 行
\begin{longtable}{p{3.6cm}cp{3cm}cp{3.6cm}}
\caption{默认值、文献范围与消融网格。}\label{tab:hyper}\\
\toprule
\textbf{超参数} & \textbf{默认} & \textbf{文献范围} & \textbf{推荐} & \textbf{消融网格} \\
\midrule\endfirsthead
\toprule
\textbf{超参数} & \textbf{默认} & \textbf{文献范围} & \textbf{推荐} & \textbf{消融网格} \\
\midrule\endhead

Token 预算 $B$ & 64 & Flamingo $\leq 64$; BLIP-2 $=32$ & 128 & $\{32,64,128\}$ \\
每句 $k$ & 8 & ContextCite top-$k$ & 8 & $\{4,8,16\}$ \\
规划预算 $B_{\mathrm{plan}}$ & $\min(32,B/4)$ & BLIP-2 32 queries & 32 & $\{0,16,32\}$ \\
伪影阈值 $\tau_A$ & 0.85 & 临床 SNR/CNR & 0.85 & $\{0.75,0.85,0.90\}$ \\
门控陡峭度 $k_{\mathrm{gate}}$ & 15 & 工程选择 & 15 & $\{10,15,20\}$ \\
权重 $(w_e,w_v,w_b)$ & $(0.5,0.3,0.2)$ & 多信号融合 & 默认 & $\pm 0.1$ 扰动 \\
评分融合 $\lambda_h$ & 0.6 & 无统一标准 & 0.6 & $\{0.4,0.6,0.8\}$ \\
边界混合 $\beta$ & 0.1 & 工程选择 & 0.1 & $\{0,0.05,0.1,0.2\}$ \\
空间先验 $\lambda_{\mathrm{spatial}}$ & 0.3 & CT-GLIP grounded & 0.3 & $\{0,0.1,0.3,0.5\}$ \\
IoU 阈值 $\tau_{\mathrm{IoU}}$ & 0.1 & 弱监督标准 & 0.04 & $\{0.03,0.04,0.05,0.1\}$ \\
InfoNCE 温度 $\tau$ & 0.07 & MoCo 0.07; SimCLR 0.1 & 0.07 & $\{0.05,0.07,0.1,0.2\}$ \\
惩罚缩放 $\gamma$ & 2.0 & 工程选择 & 2.0 & $\{1.0,2.0,3.0\}$ \\
Anatomy tiebreak $\epsilon$ & 0.05 & 工程选择 & 0.05 & $\{0.01,0.05,0.1\}$ \\
R1 ratio 阈值 $\theta_{\mathrm{ratio}}$ & 1.0 & 工程选择 & 0.6 & $\{0.4,0.6,0.8,1.0\}$ \\
R2 支持率 $\theta_{\mathrm{support}}$ & 1.0 & 工程选择 & 0.8 & $\{0.6,0.8,1.0\}$ \\
Holm 校正 & 是 & 标准 FWER & 是 & 固定 \\
Bootstrap $R$ & 5000 & 1k--10k & 5000 & $\{1000,5000,10000\}$ \\
\bottomrule
\end{longtable}


%=================================================================
\section{未决问题清单}
%=================================================================

% >>> 修改：已在代码中解决的标记 \checkmark，未解决的保留 □
\begin{enumerate}[label=$\square$]
    \item[\checkmark] \textbf{$\normop(\cdot)$ 与 $q(\cdot)$ 的形式化定义：}代码采用 min-max 归一化（$\normop$）和分位数排名（$q$），按病例/深度层独立计算。Ties 按原始 enumerate 顺序分配排名。
    \item[\checkmark] \textbf{3D NMS 规格：}轴对齐 3D IoU，可配置 NMS 阈值（默认 1.0 即不做 NMS），全局执行。
    \item[\checkmark] \textbf{侧别坐标系：}预处理中统一为 LPS 朝向（\texttt{sitk.DICOMOrient(img, "LPS")}），x 轴：$0=$ 患者右侧，$x_{\max}=$ 患者左侧。
    \item[$\square$] \textbf{$|\mathcal{P}_s|=0$ 处理：}当前实现抛出异常；尚未报告 $|\mathcal{P}_s|$ 分布。
    \item[$\square$] \textbf{$\lambda_{\mathrm{spatial}} \times \tau$ 联合敏感性：}提供网格消融，防止空间项在小 $\tau$ 下压倒语义路由。
    \item[\checkmark] \textbf{$\mathrm{sev}_i$ 归一化：}各规则严重度在配置中预设 $\in [0,1]$；R1--R5 到逐 token 的映射通过 \texttt{severity\_by\_rule} 字典实现。
    \item[\checkmark] \textbf{跨句一致性：}R6a（侧别冲突）和 R6b（存在/缺失冲突）已实现。
    \item[$\square$] \textbf{注意力代理约定：}当前代码使用 $k \cdot B$ 简化代理，尚未按公式~\eqref{eq:cache} 逐调用计算。
\end{enumerate}


%=================================================================
\section{方法--相关工作交叉引用}
%=================================================================

\begin{longtable}{p{1cm}p{3cm}p{7cm}p{2.5cm}}
\caption{方法 $\times$ 相关工作对照。}\label{tab:related}\\
\toprule
\textbf{方法} & \textbf{工作} & \textbf{对应关系} & \textbf{出处位置} \\
\midrule\endfirsthead
\toprule
\textbf{方法} & \textbf{工作} & \textbf{对应关系} & \textbf{出处位置} \\
\midrule\endhead

M1 & OctNet & 稀疏分层空间划分；叶节点 pooled 特征 & 摘要+引言 \\
M1 & TokenLearner & 少量自适应 token 降低 $O(n^2)$ 注意力代价 & Sec 3 \\
M2 & CLIP & 余弦相似+温度缩放 softmax & Sec 3.1 \\
M2 & MoCo & $\tau=0.07$ 默认；logits/$\tau$+CE 实现 & Sec 3.3 \\
M2 & CPC（InfoNCE） & 分类交叉熵解释；密度比；MI 下界 & Eq (4)--(5) \\
M2 & CT-GLIP & 3D grounded VLP；$\tau=0.07$；局部对齐 & 方法部分 \\
M3 & RECLAIM & 逐句引用；约束解码 & Sec 4 \\
M3 & LongCite/CoF & 逐句 top-$l$ 检索+引用抽取 & Sec 4.1 \\
M3/4 & ContextCite & top-$k$ 归因优于全上下文验证 & Fig 5 \\
M3/4 & CoVe & 规划$\to$核查$\to$修订结构 & 摘要+Sec 3 \\
M4 & 侧别错误文献 & 关键词/正则的侧别筛查 & 综述 \\
M5 & Bootstrap 综述 & 重采样分位数 CI & 方法部分 \\
M5 & Holm 校正 & 步降 FWER 控制 & 公式描述 \\
\bottomrule
\end{longtable}

\end{document}